\documentclass{article}

\usepackage[a4paper, left=1.5cm, right=1.5cm, top=2cm, bottom=2cm]{geometry}

\usepackage{../../../../components/mathtex}

\usepackage{hyperref}
\usepackage{fancyhdr}
\usepackage{ccicons}


% Configuration des en-têtes et pieds de page
\pagestyle{fancy}
\fancyhf{} % reset tout

\fancyhead[L]{DL3 Math-Info GA5}
\fancyhead[C]{Théorie des groupes}
\fancyhead[R]{2026-2027}

\fancyfoot[L]{Ewen Rodrigues de Oliveira\\\ccbyncsa}
\fancyfoot[R]{\thepage}

\begin{document}

\docTitle{Chapitre 5 : Produit semi-direct de groupes}

\subsection{Définition d'un produit semi-direct}

\theorem{Proposition}{}{false}{
Soient $N$ et $H$ deux groupes, notés multiplicativement, et $\varphi : H \to \mathrm{Aut}(N)$ un morphisme de groupes. On munit l'ensemble $G = N \times H$ de la loi de composition interne suivante : pour tous $n, n' \in N, h, h' \in H$
\[
(n, h) * (n', h') = (n\varphi(h)(n'), hh')
\]
Alors $(G, *)$ est un groupe.
}

\noindent{\textbf{Preuve :}
Montrons tout d'abord que la loi est associative. On a
\[
((n, h) * (n', h')) * (n'', h'') = (n\varphi(h)(n'), hh') * (n'', h'') = (n\varphi(h)(n')\varphi(hh')(n''), hh'h'').
\]
et
\[
(n, h) * ((n', h') * (n'', h'')) = (n, h) * (n'\varphi(h')(n''), h'h'') = (n\varphi(h)(n'\varphi(h')(n'')), hh'h'').
\]
Comme $\varphi(h) : N \to N$ est un morphisme de groupes, on a
\[
\varphi(h)(n'\varphi(h')(n'')) = \varphi(h)(n')\varphi(h)(\varphi(h')(n''))
\]
et comme $\varphi : H \to \mathrm{Aut}(N)$ est un morphisme de groupes, on a $\varphi(h) \circ \varphi(h') = \varphi(hh')$, ce qui donne le résultat.

\medskip
Montrons que $(e_N, e_H)$ est l'élément neutre :
\[
(n, h) * (e_N, e_H) = (n\varphi(h)(e_N), he_H) = (n, h) ; \quad (e_N, e_H) * (n, h) = (e_N\varphi(e_H)(n), e_H h) = (n, h).
\]

On vérifie que $(\varphi(h^{-1})(n^{-1}), h^{-1})$ est le symétrique de $(n, h)$ pour la loi $*$ :
\[
(n, h) * (\varphi(h^{-1})(n^{-1}), h^{-1}) = (n\varphi(h)\left(\varphi(h^{-1})(n^{-1})\right), hh^{-1}) = (nn^{-1}, e_H) = (e_N, e_H)
\]
car $\varphi(h) \circ \varphi(h^{-1}) = \varphi(hh^{-1}) = \mathrm{id}_N$. On a
\[
(\varphi(h^{-1})(n^{-1}), h^{-1}) * (n, h) = (\varphi(h^{-1})(n^{-1})\varphi(h^{-1})(n), h^{-1}h) = (\varphi(h^{-1})(n^{-1}n), h^{-1}h) = (e_N, e_H)
\]
car $\varphi(h^{-1}) : N \to N$ est un morphisme de groupes.
\hfill $\square$}

\bigskip

Définition 5.1.2. Le groupe $(G, *)$ ainsi obtenu s'appelle le **produit semi-direct** de $N$ par $H$ (relativement à $\varphi$) et est noté $N \rtimes_\varphi H$.

\theorem{Proposition}{}{false}{
Soient $N$ et $H$ deux groupes et $\varphi : H \to \mathrm{Aut}(N)$ un morphisme de groupes. Soit $G = N \rtimes_\varphi H$ le produit semi-direct de $N$ par $H$ relativement à $\varphi$. On a
\begin{enumerate}
    \item $H$ est (isomorphe à) un sous-groupe de $G$.
    \item $N$ est (isomorphe à) un sous-groupe distingué de $G$.
    \item L'application $G \to H$ qui à $(n, h)$ associe $h$ est un morphisme de groupes de noyau $N$ induisant un isomorphe $G/N \to H$.
    \item $G = NH = HN$.
\end{enumerate}
}

\noindent{\textbf{Preuve :}
On considère les deux applications $H \to G$ et $N \to G$ définies par $h \mapsto (e_N, h)$ et $n \mapsto (n, e_H)$. On vérifie aisément que ces deux applications sont des morphismes de groupes, et qu'ils sont injectifs. Si on note $H'$ et $N'$ leur image respective, alors $H \cong H'$ et $N \cong N'$. Dans les énoncés on identifie $H$ à $H'$ et $N$ à $N'$. La projection $G \to H$ qui à $(n, h)$ associe $h$ est un morphisme de groupes, surjectif, de noyau $N'$. Donc $N'$ est distingué dans $G$ et cette projection induit un morphisme $G/N' \to H$. Enfin comme $N'$ est distingué dans $G$, on a $N'H'$ est un sous-groupe de $G$ et $N'H' = H'N'$ (voir Proposition 2.4.7). Par ailleurs tout élément $(n, h)$ de $G = N \rtimes_\varphi H$ se décompose en $(n, h) = (n, e_H) * (e_N, h) \in N'H'$ donc $G = N'H'$.
\hfill $\square$}

\subsection{Suites exactes courtes de groupes}

Définition 5.2.1. Une **suite exacte courte de groupes** est la donnée de 3 groupes $N, G, H$, de deux morphismes de groupes $i : N \to G$ et $\pi : G \to H$ tels que $i$ soit injective, $\pi$ surjective et $i(N) = \mathrm{ker}\,\pi$. On note une telle suite comme ceci :
\[
1 \longrightarrow N \overset{i}{\longrightarrow} G \overset{\pi}{\longrightarrow} H \longrightarrow 1
\]
Une suite exacte courte de groupes est dite **scindée** s'il existe un morphisme de groupes $s : H \to G$ tel que $\pi \circ s = \mathrm{id}_H$.

\example{
\begin{itemize}
    \item Les deux suites exactes courtes suivantes sont scindées
    \[
    1 \longrightarrow \mathcal{A}_n \longrightarrow \mathcal{S}_n \overset{\epsilon}{\longrightarrow} \{\pm 1\} \longrightarrow 1
    \]
    \[
    1 \longrightarrow \mathrm{SL}_n(\mathbb{R}) \longrightarrow \mathrm{GL}_n(\mathbb{R}) \overset{\mathrm{det}}{\longrightarrow} \mathbb{R}^* \longrightarrow 1
    \]
    Pour la première on peut prendre $s : \{\pm 1\} \to \mathcal{S}_n$ l'application qui à $-1$ associe la transposition $(12)$. Pour la deuxième suite exacte, on peut prendre $s : \mathbb{R}^* \to \mathrm{GL}_n(\mathbb{R})$ l'application qui à $\lambda$ associe la matrice diagonale $\mathrm{Diag}(\lambda, 1, \dots, 1)$. On vérifie que ces deux sections sont des morphismes de groupes.
    
    \item Pour $\varphi : H \to \mathrm{Aut}(N)$ on a une suite exacte courte scindée
    \[
    1 \longrightarrow N \longrightarrow N \rtimes_\varphi H \longrightarrow H \longrightarrow 1
    \]
    
    \item La suite exacte courte suivante n'est pas scindée
    \[
    1 \longrightarrow \mathbb{Z} \overset{\times 2}{\longrightarrow} \mathbb{Z} \overset{\pi}{\longrightarrow} \mathbb{Z}/2\mathbb{Z} \longrightarrow 1
    \]
    car si elle était scindée, on aurait un morphisme de groupes non nul $s : \mathbb{Z}/2\mathbb{Z} \to \mathbb{Z}$, ce qui est absurde car $\mathbb{Z}$ ne contient pas d'élément d'ordre 2.
\end{itemize}
}

\theorem{Théorème}{}{false}{
Soit
\[
1 \longrightarrow N \overset{i}{\longrightarrow} G \overset{\pi}{\longrightarrow} H \longrightarrow 1
\]
une suite exacte courte scindée de groupes. Alors il existe un morphisme de groupes $\varphi : H \to \mathrm{Aut}(N)$ et un morphisme de groupes $f : N \rtimes_\varphi H \to G$ tel que $\pi \circ f$ soit la projection sur $H$.
}

\noindent{\textbf{Preuve :}
On identifie $N$ à $i(N) = \mathrm{ker}\,\pi$ et on note $s : H \to G$ une section de $\pi$. On remarque que pour tout $h \in H$, $s(h)n s(h^{-1}) \in N$. En effet
\[
\pi(s(h)n s(h^{-1})) = \pi(s(h))\pi(n)\pi(s(h^{-1})) = h e_H h^{-1} = e_H.
\]
Par ailleurs $n \mapsto s(h)n s(h^{-1})$ est un automorphisme du groupe $N$ que l'on note $\varphi(h)$. Comme $s$ est un morphisme de groupes, on a $\varphi(h) \circ \varphi(h') = \varphi(hh')$ donc $\varphi : H \to \mathrm{Aut}(N)$ est un morphisme de groupes. Montrons que l'application $f : N \rtimes_\varphi H \to G$ définie par $f(n, h) = n s(h)$ est un morphisme de groupes bijectif. La suite d'égalités
\[
f((n, h)*(n', h')) = f(n\varphi(h)(n'), hh') = f(n s(h)n' s(h)^{-1}, hh') = n s(h)n' s(h') = f(n, h)f(n', h')
\]
montre que $f$ est un morphisme de groupes. Supposons $n s(h) = e_G$. En appliquant $\pi$ on obtient $\pi(n s(h)) = \pi(n)\pi(s(h)) = h = e_H$ et $n = e_G = e_N$. Ainsi $f$ est injective. Fixons $g \in G$. On pose $n = g s(\pi(g^{-1}))$ et on vérifie que $n \in N$: $\pi(n) = \pi(g)\pi(s(\pi(g^{-1}))) = \pi(g)\pi(g^{-1}) = e_H$. Ainsi $g = n s(\pi(g)) = f(n, \pi(g))$ et $f$ est surjective.
\hfill $\square$}


\subsection{Exemples}

\subsubsection{Le groupe $\mathcal{S}_3$}

Le sous-groupe de $\mathcal{S}_3$ engendré par le cycle $(123)$ est distingué dans $\mathcal{S}_3$, on le note $N$. On considère le sous-groupe $H$ d'ordre 2 engendré par $(12)$.

\theorem{Proposition}{}{false}{
Le groupe symétrique $\mathcal{S}_3$ est isomorphe à $N \rtimes_\varphi H$ avec $\varphi : H \to \mathrm{Aut}(N)$ défini par $\varphi(\mathrm{id}) = \mathrm{id}$ et $\varphi((12))(x) = (12)x(12)$.
}

\noindent{\textbf{Preuve :}
Soit $\psi : N \rtimes_\varphi H \to G$ l'application qui à $(x, y)$ associe $xy$. Il faut montrer que $\psi$ est un morphisme de groupes. Si tel est le cas, il sera bijectif car $\psi$ est surjectif.

On a
\[
\psi((n, h) * (n', h')) = \psi(n\varphi(h)(n'), hh') = n\varphi(h)(n')hh'
\]
Montrons que pour tout $n' \in N, h \in H$ on a
\[
\varphi(h)(n')h = hn'.
\]
Si $h = \mathrm{id}$ c'est vérifié. Si $h = (12)$, calculons
\[
\varphi(h)(n')h = (12)n'(12)(12) = (12)n' = hn'.
\]
D'où l'isomorphisme.
\hfill $\square$}


\subsubsection{Le groupe diédral}

Le groupe diédral d'ordre $2n$ noté $D_n$ peut être défini comme le groupe des isométries d'un polygone régulier à $n$ côtés. En fixant 0 l'isobarycentre d'un tel polygone et en fixant un sommet $A_0$ de ce polygone, on définit un repère orthonormé $(0, \vec{i}, \vec{j})$ de $\mathbb{R}^2$ ce qui nous permet d'identifier $\mathbb{R}^2$ muni de ce repère à $\mathbb{C}$ via la bijection $x\vec{i} + y\vec{j} \mapsto x + iy$. Notons $\sigma$ la rotation d'angle $\frac{2\pi}{n}$ et $\tau$ la symétrie d'axe $(0x)$. L'expression dans $\mathbb{C}$ de $\sigma$ et $\tau$ sont données par
\[
\sigma(z) = e^{\frac{2i\pi}{n}}z, \quad \tau(z) = \bar{z}.
\]

On sait que le groupe des isométries vectorielles du plan est l'ensemble des rotations et des symétries (c'est $O_2(\mathbb{R})$). Le groupe des isométries d'un polygone régulier à $n$ côtés, avec l'identification précédente, est le sous-groupe de $O_2(\mathbb{R})$ engendré par $\sigma$ et $\tau$. Remarquons que $\sigma$ est d'ordre $n$ et $\tau$ est d'ordre 2. De plus $\tau\sigma\tau = \sigma^{-1}$, donc
\[
D_n = \langle \sigma \rangle \cup \langle \sigma \rangle \tau
\]
a exactement $2n$ éléments. En notant $H = \langle \tau \rangle \cong \mathbb{Z}/2\mathbb{Z}$ et $N = \langle \sigma \rangle \cong \mathbb{Z}/n\mathbb{Z}$ on a $\varphi : H \to \mathrm{Aut}(N)$ donnée par $\varphi(\tau)(\sigma) = \tau\sigma\tau^{-1} = \sigma^{-1}$ est un morphisme de groupes et
\[
D_n \cong \mathbb{Z}/n\mathbb{Z} \rtimes_\varphi \mathbb{Z}/2\mathbb{Z}.
\]
\newpage
\tableofcontents
\section*{Crédits}
Ce cours provient du polycopié de cours de l'Université Paris Cité de Muriel Livernet enrichi et modifié par mes soins à des fins pédagogiques. Assurez-vous de citer correctement la source si vous utilisez ce cours dans vos travaux.\\\\
Merci à Laurent Cai à son aide dans la rédaction des démonstrations.\\\\
Ce cours est mis à disposition selon les termes de la Licence Creative Commons \ccbyncsa\ \textbf{Attribution - Pas d'Utilisation Commerciale - Partage dans les Mêmes Conditions 4.0 International}. 
Pour voir une copie de cette licence, visitez \url{http://creativecommons.org/licenses/by-nc-sa/4.0/}.

\end{document}